
\subsubsection{3.1 Overview}

The methodology comprises three components: (1) model collection from HuggingFace Hub, (2) heavy-tailed exponent computation via the Hill estimator, and (3) variance analysis across model populations.

WeightWatcher, the canonical implementation of HT-SR analysis by Martin and Mahoney, was used to ensure methodological continuity with prior CNN analysis while testing on a new architecture class.

\subsubsection{3.2 Model Collection}

Vision Transformer models were collected from HuggingFace Model Hub using the following criteria:

\begin{itemize}
\item Filter: pipeline\_tag = image-classification
\item Search: model name contains ''vit''
\item Sort: downloads (descending)
\item Target: 100 models
\end{itemize}

Filtering by downloads prioritizes well-established models with community validation, reducing noise from experimental checkpoints.

\subsubsection{3.3 Heavy-Tailed Exponent Computation}

For each model, heavy-tailed exponent $\alpha$ was computed using WeightWatcher. WeightWatcher performs singular value decomposition on each weight matrix and fits a power-law distribution $p(x) \propto x^{-\alpha}$ using the Hill estimator. The $\alpha$ value characterizes tail behavior: lower values indicate heavier tails.

\subsubsection{3.4 Variance Analysis}

Two variance metrics were computed:

\begin{enumerate}
\item \textbf{Global variance} $\sigma _global$: Standard deviation of mean $\alpha$ across all models
\item \textbf{Within-family variance} $\sigma _family$: Standard deviation within model families sharing training origin
\end{enumerate}

The gate condition was set at $\sigma$ < 0.5 for bounded variance.

